\DocumentMetadata{
  pdfversion=2.0,
  pdfstandard=ua-2,
  lang=en-US,
  tagging=on,
  tagging-setup={math/setup=mathml-SE,tabsorder=structure}
}
\documentclass[11pt,letterpaper]{article}
\usepackage[margin=0.68in,includefoot]{geometry}
\usepackage{newcomputermodern}
\usepackage{amsmath,amscd,mathtools}
\usepackage{tikz,adjustbox}
\providecommand{\square}{\mdlgwhtsquare}
\usepackage[hidelinks]{hyperref}
\hypersetup{
  pdftitle={Numerical Analysis PhD Exam, January 1996},
  pdfauthor={University of Florida Department of Mathematics},
  pdfsubject={Graduate mathematics examination},
  pdfdisplaydoctitle=true,
  bookmarksopen=true
}
\setcounter{secnumdepth}{0}
\setlength{\parindent}{0pt}
\setlength{\parskip}{0.28em}
\setlength{\emergencystretch}{3em}
\setlength{\leftmargini}{2.15em}
\setlength{\leftmarginii}{2.65em}
\setlength{\leftmarginiii}{3em}
\renewcommand{\labelenumi}{\textbf{\theenumi.}}
\pagestyle{empty}
\raggedbottom
\allowdisplaybreaks
\newenvironment{examproblems}[1]{%
  \begin{enumerate}%
  \setcounter{enumi}{#1}%
  \setlength{\topsep}{0.35em}%
  \setlength{\partopsep}{0pt}%
  \setlength{\itemsep}{0.48em}%
  \setlength{\parsep}{0.18em}%
}{\end{enumerate}}
\newenvironment{examsubparts}{%
  \begin{enumerate}%
  \setlength{\topsep}{0.18em}%
  \setlength{\partopsep}{0pt}%
  \setlength{\itemsep}{0.16em}%
  \setlength{\parsep}{0.1em}%
}{\end{enumerate}}
\newenvironment{examinstructions}{%
  \par\begingroup\itshape
}{%
  \par\endgroup\vspace{0.18em}
}
\makeatletter
\renewcommand\section{\@startsection{section}{1}{0pt}%
  {-1.25ex plus -.25ex minus -.1ex}%
  {0.55ex plus .1ex}%
  {\normalfont\large\bfseries}}
\renewcommand{\@maketitle}{%
  \newpage\null\vskip -1em%
  {\footnotesize\bfseries
    \href{https://www.ufl.edu/}{University of Florida}, \href{https://math.ufl.edu/}{Department of Mathematics}\par}%
  \vskip -0.5em%
  \tagstructbegin{tag=Title}%
    {\LARGE\bfseries\href{https://gma.math.ufl.edu/exams/numerical-analysis-phd/}{Numerical Analysis PhD Exam}, January 1996\par}%
  \tagstructend%
  \vskip 0.25em%
  \tagmcbegin{artifact}%
    \hrule height 1pt%
  \tagmcend%
  \vskip 1em%
}
\makeatother
\title{Numerical Analysis PhD Exam, January 1996}
\author{}
\date{}
\begin{document}
\maketitle
\thispagestyle{empty}
\begin{examproblems}{0}
\item
Let~\(A\) be a symmetric, positive definite matrix. Show that if~\(A\) is \(LU\) factored, then the elements of~\(U\) in absolute value are bounded by the largest diagonal element of~\(A\).
\item
Suppose the columns of a matrix~\(A\) are independent. Given the \(QR\) factorization of~\(A\), give a formula for the distance from some given vector~\(b\) to the space spanned by the columns of~\(A\).
\item
Give a careful statement and proof of Gershgorin’s Theorem. Find the upper and lower bounds for the set of eigenvalues for the \(4\times 4\) Hilbert matrix having \((i,j)\)-th entry \((i+j-1)^{-1}\).
\item
Explain how the singular values of a matrix can be used to determine the rank of a matrix in the presence of rounding errors.
\item
Define the term “algorithm”. Explain what is meant by saying that an algorithm is well conditioned and numerically stable. Carefully formulate a numerically stable algorithm for computing the positive solution of the quadratic equation \(x^2+2px-q=0\), where~\(p\) and~\(q\) are arbitrary positive numbers. Justify the numerical stability of your algorithm.
\item
Determine the cubic spline \(s(x)\) on the grid \(\{-1,0,1\}\) with \(s(-1)=s(0)=0\), \(s(1)=4\), and \(s''(-1)=s''(1)=0\).
\item
For \(h>0\), let \(T(h)\) be the trapezoidal approximation of \(\tau=\int_a^b f(x)\,dx\), using an equi-spaced grid of step size~\(h\).
\begin{examsubparts}
\item[\textnormal{(a)}]
What is the order of the error in this approximation as \(h\to 0\)?
\item[\textnormal{(b)}]
Assuming that \(T(h)\) can be expressed as a power series of even powers of~\(h\), show how a more accurate approximation can be obtained by using \(T(h)\) and \(T(h/2)\). Express this approximation in terms of \(T(h),T(h/2)\). (That is, you need to verify the first step in Romberg integration.)
\item[\textnormal{(c)}]
Describe the method of Romberg integration. Formulate it in terms of a table (Neville type).
\end{examsubparts}
\item
Let~\(A\) be a nonsingular \(n\times n\) matrix and \(\{X_k\}\) a sequence of \(n\times n\) matrices generated by the iteration: 
\[
X_{k+1}=X_k+X_k(I-AX_k).
\]
 Show that if \(X_0\) is chosen such that the spectral radius of \(I-AX_0\) is strictly less than 1, then \(\{X_k\}\) converges to the inverse of~\(A\).
\item
State Euler’s method for approximating the solution to the differential equation 
\[
\frac{dx}{dt}=f(x(t)),\qquad x(0)=x_0.
\]
 Show that the error in Euler’s method is \(O(\Delta t)\), as \(\Delta t\to 0\), for~\(t\) sufficiently small, where \(\Delta t\) is the step size.
\end{examproblems}
\end{document}
